%%%%%%%%%%%%%%%%%%%%%%%%%%%%%%%%%%%%%%%%%
% Academic CV - Jonathan Lozano
% Based on the ModernCV template by Xavier Danaux
% License: CC BY-NC-SA 3.0
%%%%%%%%%%%%%%%%%%%%%%%%%%%%%%%%%%%%%%%%%

%----------------------------------------------------------------------------------------
%   PACKAGES AND DOCUMENT CONFIGURATION
%----------------------------------------------------------------------------------------

\documentclass[11pt,a4paper,roman]{moderncv}

\moderncvstyle[nosymbols]{classic}   % 'casual', 'classic', 'banking', 'oldstyle', 'fancy'
\moderncvcolor{blue}      % 'blue', 'orange', 'green', 'red', 'purple', 'grey', 'black'

\usepackage[scale=0.84]{geometry}
\setlength{\hintscolumnwidth}{2.3cm} % Width of the dates column
\usepackage[T1]{fontenc}
\usepackage{mathpazo} % Palatino text with matched math
\linespread{1.05}     % Palatino benefits from slightly more leading
\microtypesetup{expansion=false} % marvosym symbol font is not scalable

% Restore the marvosym envelope/house icons in the header.
% (\marvosymbol is provided by the letters collection loaded via [nosymbols];
%  codes 66 and 205 are the marvosym collection's email and homepage glyphs.)
\renewcommand*{\emailsymbol}{\marvosymbol{66}~}
\renewcommand*{\homepagesymbol}{{\Large\marvosymbol{205}}~}

% Tighten the spacing produced by \cventry / \cvitem
\setlength{\parskip}{0pt}
\renewcommand*{\quotefont}{\normalsize\normalfont}

%----------------------------------------------------------------------------------------
%   PERSONAL INFORMATION
%----------------------------------------------------------------------------------------

\firstname{Jonathan}
\familyname{Lozano}

\title{PhD Candidate in Theoretical Physics}
\address{Weinberg Institute for Theoretical Physics}{The University of Texas at Austin}{}
\email{jonathanloz@utexas.edu}
\homepage{jonathanloz.github.io}
\extrainfo{ORCID: \href{https://orcid.org/0000-0002-9638-5173}{0000-0002-9638-5173}}

%----------------------------------------------------------------------------------------

\begin{document}

\makecvtitle

%----------------------------------------------------------------------------------------
%   EDUCATION
%----------------------------------------------------------------------------------------

\section{Education}

\cventry{2023--present}{Ph.D.\ in Physics}{The University of Texas at Austin}{USA}{}{%
Advisor: Dr.\ Antonios Alvertis. \newline
GPA: 3.9/4.00. Research in condensed matter and many-particle physics.}

\cventry{2016--2021}{B.Sc.\ in Physics}{National Autonomous University of Mexico (UNAM)}{Mexico City}{\textit{GPA: 9.39/10.0 --- High Honors}}{%
Advisor: Dr.\ Manuel Torres Labansat. Coursework included two graduate-level
courses in quantum field theory and a graduate course in differential geometry
and topology for physicists.\newline
\emph{Thesis:} Spontaneous symmetry breaking and extended field configurations
in a scalar theory subject to a potential with two families of vacuum states.}

%----------------------------------------------------------------------------------------
%   PUBLICATIONS
%----------------------------------------------------------------------------------------

\section{Publications}

\cvitem{2026}{\textbf{J. Lozano-Mayo}. Generalized virial identities: radial
constraints for solitons, instantons, and bounces. \textit{J. High Energ.
Phys.} \textbf{04} (2026) 058.
\href{https://doi.org/10.1007/JHEP04(2026)058}{doi:10.1007/JHEP04(2026)058}}
\cvitem{2026}{\textbf{J. Lozano-Mayo} and M. Torres-Labansat. Multi-kinks and
composite oscillons in a commensurable and non degenerate double sine-Gordon
model. \textit{J. High Energ. Phys.} \textbf{09} (2026) 218.
\href{https://doi.org/10.1007/JHEP09(2026)218}{doi:10.1007/JHEP09(2026)218}}
\cvitem{2025}{J. J. Ziegler, K. Freese, \textbf{J. Lozano}, G. Montefalcone.
Explaining the ``too massive'' high-redshift galaxies in JWST data: a numerical
study of three effects and a simple relation. \textit{(In revision).}}

\cvitem{2021}{\textbf{J. Lozano-Mayo} and M. Torres-Labansat. Kink solutions in
a generalized scalar $\phi_G^4$ field model. \textit{J. Phys. Commun.}
\textbf{5} (2021) 025004.
\href{https://doi.org/10.1088/2399-6528/abdd83}{doi:10.1088/2399-6528/abdd83}}

\cvitem{2024-2025}{
\textit{ALICE Collaboration}: Contributed to ITS performance analysis, and hadron-$D^0$ azimuthal-correlation studies as service work during December 2024-October 2025. Publications bearing my name under ALICE's collective authorship policy are not listed here.}

%----------------------------------------------------------------------------------------
%   TALKS
%----------------------------------------------------------------------------------------

\section{Conference Presentations}

\cvitem{2025}{\textbf{J. L. Mayo}. \emph{$h$--$D^0$ angular correlations and
$\Delta$ROF performance studies at ITS2}. ALICE-USA Meeting.}

\cvitem{2021}{\textbf{J. L. Mayo}. \emph{Multi-solitones en teor\'ias escalares
de campo con vac\'ios no-degenerados: el modelo doble de sine-Gordon}. LXVI
National Physics Congress, Mexican Physical Society.}

%----------------------------------------------------------------------------------------
%   RESEARCH EXPERIENCE
%----------------------------------------------------------------------------------------

\section{Research Experience}
\cventry{2026}{Strongly interacting electron-phonon systems}{UT Austin}{}{Advisor: Prof.\ Antonios Alvertis}{%
Studied the properties of strongly interacting electron-phonon systems using tensor network methods.}
\cventry{2025}{$h$--$D^0$ Angular Correlations}{ALICE Collaboration}{}{Advisor: Prof.\ Christina Markert}{%
Studied the properties of charm production across distinct color environments
using ALICE Run-3 Pb--Pb collisions, characterizing the azimuthal angular
difference between a high-momentum trigger hadron and a lower-momentum
associated $D^0$ meson.}

\cventry{2024}{JWST Early Massive Galaxies and Top-Heavy IMFs}{UT Austin}{}{Advisor: Prof.\ Katherine Freese}{%
Modeled the spectral energy distributions of high-redshift galaxies with the
population-synthesis code P\'EGASE to constrain the parameter space in which
JWST observations remain consistent with the standard $\Lambda$CDM cosmology.}

\cventry{2021--2022}{Multi-kinks in Scalar Field Theories with Non-degenerate Vacua}{UNAM}{}{Advisor: Dr.\ Manuel Torres Labansat}{%
Investigated the formation of static $n$-kink structures in models with
deformed potentials. Performed numerical simulations in Julia to characterize
the binding forces between kinks and derived an analytical relation for the
energy of the static multi-kink.}

\cventry{2020--2021}{Quantum Renormalization of Kink Masses and Virial Relations}{UNAM}{}{Advisor: Dr.\ Manuel Torres Labansat}{%
Renormalized the one-loop quantum mass correction of kink configurations
arising from the topology of a generalized $\phi^4$ potential using
perturbative techniques. Established the existence of static multi-kink
configurations via asymptotic analysis and developed a framework for their
stability.}

\cventry{2018--2019}{Photon Wave Function}{UNAM}{}{Advisor: Dr.\ Manuel Torres Labansat}{%
Examined whether a well-defined photon wave function can be constructed from
the photon dispersion relation. Used group-theoretic methods to build the
quantum operators and Lorentz transformations acting on a proposed six-component
wave function, then derived the Lagrangian density and verified the expected
symmetries of the theory.}

\cventry{2019}{Time-of-Flight Mass Spectrometry}{UNAM}{}{Advisor: Dr.\ Juan L\'opez Pati\~no}{%
Produced proton--air collisions with a low-energy linear collider and analyzed
the reaction products using the time-of-flight mass spectrometry technique.}

%----------------------------------------------------------------------------------------
%   TEACHING & ACADEMIC SERVICE
%----------------------------------------------------------------------------------------

\section{Teaching \& Academic Service}

\cventry{2026}{Teaching Assistant --- Graduate Quantum Mechanics}{UT Austin}{}{}{%
Mentored students and graded assignments for a graduate course on Quantum Mechanics.}
\cventry{2025}{Lecturer --- Quantum Field Theory I}{ICTP (International Centre for Theoretical Physics)}{}{}{%
Delivered an introductory quantum field theory course as part of the
International Centre for Theoretical Physics outreach initiative.}

\cventry{2024}{Teaching Assistant --- Advanced Particle Physics}{ICTP}{}{}{%
Led problem-set discussions and graded assignments for a graduate course on
advanced topics in quantum field theory.}
\cventry{2025}{Teaching Assistant --- Waves and Optics}{UT Austin}{}{}{%
Held weekly discussion sections and graded coursework spanning wave phenomena, optics, and electromagnetism.}
\cventry{2024--2025}{Teaching Assistant --- Modern Physics}{UT Austin}{}{}{%
Led weekly discussion sections and graded coursework spanning statistical
mechanics, quantum mechanics, and special relativity.}

\cventry{2023}{Teaching Assistant --- Electromagnetism Laboratory (PHY 105N)}{UT Austin}{}{}{%
Led two weekly laboratory sections and supported the experimental setups.}

\cventry{2022}{Teaching Assistant --- Thermodynamics}{UNAM}{}{}{%
Held problem-solving sessions and created and graded homework and exams.}

\cventry{2020}{Teaching Assistant --- Nuclear and Subnuclear Physics}{UNAM}{}{}{%
Mentored students and graded homework and exams for a course taught by Prof.\
Manuel Torres Labansat.}

\cventry{2022}{Referee}{J.\ Phys.\ G: Nuclear and Particle Physics}{}{}{%
Reviewed submitted manuscripts.}

\cventry{2017}{Volunteer Instructor}{UNAM School of Humanities and Sciences}{}{}{%
Taught physics and mathematics in an undergraduate preparation course.}

%----------------------------------------------------------------------------------------
%   AWARDS
%----------------------------------------------------------------------------------------

\section{Awards \& Distinctions}

\cvitem{2022}{\textbf{DPG/IAPS PLANCKS Munich Travel Grant} --- awarded to attend the PLANCKS competition at LMU Munich.}
\cvitem{2022}{\textbf{Winner, PLANCKS 2022 Mexican Preliminary} --- placed 1st of 30 teams in the Mexican Physics Tournament (\emph{Torneo Mexicano de F\'isica}).}
\cvitem{2021}{\textbf{Juan Manuel Lozano Mej\'ia Diploma} --- awarded by the UNAM Institute of Physics to students with the most outstanding thesis research.}
\cvitem{2021}{\textbf{Excellence Cluster Fellowship} --- PRISMA, Johannes Gutenberg University Mainz.}
\cvitem{2021}{\textbf{Top 10, PLANCKS 2021 (Porto)} --- placed 7th of 50 teams in the international theoretical physics competition.}
\cvitem{2021}{\textbf{Winner, PLANCKS 2021 Mexican Preliminary} --- selected for the first Mexican team to attend PLANCKS after placing 1st nationally.}
\cvitem{2021}{\textbf{Honorific Mention}, B.Sc.\ thesis defense (UNAM) --- highest distinction awarded at defense.}
\cvitem{2019}{\textbf{Research Assistantship}, UNAM Institute of Physics.}
\cvitem{2018}{\textbf{AMC Scholarship} --- Mexican Academy of Sciences, XXVIII Scientific Research Summer.}
\cvitem{2018}{\textbf{ICF-UNAM Scholarship} --- VII Experimental Physics School.}
\cvitem{2017}{\textbf{IF-UNAM Scholarship} --- XXV Physics School.}

%----------------------------------------------------------------------------------------
%   ADDITIONAL TRAINING
%----------------------------------------------------------------------------------------

\section{Additional Training}

\cvitem{2025}{Summer School on CMB Data Analysis --- UT Austin.}
\cvitem{2022}{Summer School on Modeling and Tools for Data Analysis in Science --- UANL.}
\cvitem{2021}{Bad Honnef Physics School (DPG) --- Methods of Effective Field Theory and Lattice Field Theory.}
\cvitem{2021}{VII Mexican School on String Theory and Supersymmetry (MSSS) --- University of Guanajuato.}
\cvitem{2021}{Mexican Astro-Particle School (MAPS) --- University of Guanajuato.}

%----------------------------------------------------------------------------------------
%   RESEARCH INTERESTS
%----------------------------------------------------------------------------------------

\section{Research Interests}

\cvitem{}{Condensed matter physics \quad$\cdot$\quad Many-particle physics \quad$\cdot$\quad Quantum phenomena
\quad$\cdot$\quad Topological and
non-topological solitons \quad$\cdot$\quad Heavy-flavor physics}

%----------------------------------------------------------------------------------------
%   SKILLS
%----------------------------------------------------------------------------------------

\section{Technical Skills}

\cvitem{Programming}{Julia, Python, C++, Mathematica, \LaTeX, HTML}
\cvitem{Scientific tools}{ROOT/O$^2$ (ALICE framework), QtiPlot, Origin, Tracker, ImageJ}

\section{Languages}

\cvitem{}{\textbf{Spanish} (native) \quad$\cdot$\quad \textbf{English} (C1)
\quad$\cdot$\quad \textbf{German} (A2) \quad$\cdot$\quad \textbf{French} (A2)}

\end{document}
